\section{Sello dodecafásico y lectura racional reversible}
\label{sec:sello}

El registro terminal tiene una coordenada visible y una coordenada
firmada. El sello se obtiene de esta última mediante la transformación
de Hadamard por órbitas; una lectura escalar posterior permite recuperar
el sello cuando se conserva su marco de representación. Esta propiedad
precisa qué significa aquí una definición estructural de una constante:
se explicitan los mapas de producción y recuperación, además de su valor.

Sean $U\in\Z^{12}$ la coordenada firmada producida por el registro,
$H_4$ la matriz de Hadamard desarrollada en la sección~\ref{sec:action}
y $P_{\rm orb}$ la permutación que agrupa
\[
 (1,4,7,10),\qquad(2,5,8,11),\qquad(3,6,9,12).
\]
Defínase
\[
 \mathcal H_{12}=P_{\rm orb}^{-1}
 (H_4\oplus H_4\oplus H_4)P_{\rm orb}.
\]
Se cumple $\mathcal H_{12}^2=4I$. Cuando
$\mathcal H_{12}U\equiv0\pmod4$, el sello integral es
\[
 K=\tfrac14\mathcal H_{12}U,
 \qquad U=\mathcal H_{12}K.
\]
Las diferencias transversales $D_3K,D_4K$ y la carga $Q(K)$ recuperan
los canales especificados por el registro \cite{registro}.
El sello es así una coordenada reconstruible de ese registro orientado;
la profundidad y la historia completa siguen siendo campos distintos.

\begin{proposition}[Recuperación de bloques y transporte del origen]
Fijadas una base entera $b\ge2$ y doce posiciones ordenadas, sea
$K_j\in\{0,\ldots,b-1\}$. Defínanse
\[
 I(K)=\sum_{j=1}^{12}K_jb^{12-j},\qquad
 \kappa(K)=\frac{I(K)}{b^{12}-1}.
\]
La lectura exacta $\kappa$ determina unívocamente los doce bloques.
Para la rotación $\rho K=(K_2,\ldots,K_{12},K_1)$,
\[
 \boxed{\kappa(\rho K)=b\kappa(K)-K_1.}
\]
\end{proposition}
\begin{proof}
El entero $I=(b^{12}-1)\kappa$ se descompone de manera única en doce
posiciones de base $b$, rellenando con ceros las posiciones iniciales
necesarias. Por otra parte,
\[
 I(\rho K)=bI(K)-K_1(b^{12}-1).
\]
Dividir por $b^{12}-1$ da la ley de transformación. Para el bloque
constante $K_j=b-1$, la lectura es $\kappa=1$; el tipo de representación
periódica fijado conserva igualmente su recuperación.
\end{proof}

En la base $1000$ del registro, los bloques de tres cifras permanecen
recuperables, incluidos los ceros iniciales. La lectura racional
periódica conserva más que una aproximación decimal truncada.
Al modificar el origen de fase, cambia generalmente su valor según
la identidad anterior; al completar doce rotaciones, retorna.
El avance de memoria del estado TPK es una transformación adicional:
la periodicidad del lector no identifica las historias acumuladas.

\begin{apunte}{lectura reversible e información}
La recuperación exacta de los bloques es una propiedad del lector
racional en su marco ordenado. La complejidad de palabras, la entropía
de una dinámica simbólica, la entropía del estado reducido y la
producción termodinámica son objetos distintos. Una relación entre
ellos requiere especificar la distribución o el estado y el mapa de
transporte correspondiente; la reversibilidad del sello, por sí sola,
no identifica esos funcionales.
\end{apunte}
